\documentclass[a4paper,12pt]{article}

% Русский язык и кодировка
\usepackage[utf8]{inputenc}
\usepackage[T2A]{fontenc}
\usepackage[russian]{babel}

% Математика
\usepackage{amsmath}
\usepackage{amssymb}
\usepackage{mathtools}

\begin{document}
	
	\section*{Решение системы}
	
	Рассмотрим систему
	
	\begin{equation}
		\beta
		\begin{pmatrix}
			0\\
			1\\
			2x\\
			3x^2
		\end{pmatrix}
		=
		-\omega_1
		\begin{pmatrix}
			1\\
			-1\\
			1\\
			-1
		\end{pmatrix}
		+
		\omega_2
		\begin{pmatrix}
			1\\
			t_1\\
			t_1^2\\
			t_1^3
		\end{pmatrix}
		-
		(1-\omega_1-\omega_2)
		\begin{pmatrix}
			1\\
			t_2\\
			t_2^2\\
			t_2^3
		\end{pmatrix},
	\end{equation}
	
	а также
	
	\begin{equation}
		\frac{t_3-t_2}{t_1+1}=1,
	\end{equation}
	
	\begin{equation}
		\frac{t_3-t_2}{t_2-t_1}
		=
		\frac{t_3-t_2}{t_1+1}.
	\end{equation}
	
	Необходимо найти
	
	\[
	(\omega_1,\omega_2,t_1,t_2,t_3,\beta).
	\]
	
	Будем считать $x$ заданным параметром.
	
	\subsection*{1. Дополнительные уравнения для $t_1,t_2,t_3$}
	
	Из первого дополнительного уравнения
	
	\[
	\frac{t_3-t_2}{t_1+1}=1
	\]
	
	получаем
	
	\begin{equation}
		t_3-t_2=t_1+1.
		\label{eq:t3t2}
	\end{equation}
	
	Из второго дополнительного уравнения
	
	\[
	\frac{t_3-t_2}{t_2-t_1}
	=
	\frac{t_3-t_2}{t_1+1}
	\]
	
	и уже полученного равенства
	
	\[
	\frac{t_3-t_2}{t_1+1}=1
	\]
	
	следует
	
	\[
	\frac{t_3-t_2}{t_2-t_1}=1.
	\]
	
	Следовательно,
	
	\begin{equation}
		t_3-t_2=t_2-t_1.
		\label{eq:equal}
	\end{equation}
	
	Таким образом, промежутки между $t_1,t_2,t_3$ одинаковы:
	
	\[
	t_2-t_1=t_3-t_2.
	\]
	
	Из \eqref{eq:t3t2} имеем
	
	\[
	t_2-t_1=t_1+1,
	\]
	
	откуда
	
	\begin{equation}
		t_2=2t_1+1.
		\label{eq:t2}
	\end{equation}
	
	А тогда
	
	\[
	t_3=t_2+t_1+1,
	\]
	
	то есть
	
	\begin{equation}
		t_3=3t_1+2.
		\label{eq:t3}
	\end{equation}
	
	Таким образом,
	
	\begin{equation}
		\boxed{
			t_2=2t_1+1,
			\qquad
			t_3=3t_1+2.
		}
	\end{equation}
	
	При этом из исходного дополнительного уравнения необходимо
	
	\begin{equation}
		t_1\neq -1.
	\end{equation}
	
	\subsection*{2. Разбиение векторного уравнения на скалярные}
	
	Из векторного равенства получаем четыре уравнения:
	
	\begin{align}
		0
		&=
		-\omega_1+\omega_2-(1-\omega_1-\omega_2),
		\label{eq:1}\\
		\beta
		&=
		\omega_1+\omega_2t_1
		-(1-\omega_1-\omega_2)t_2,
		\label{eq:2}\\
		2\beta x
		&=
		-\omega_1+\omega_2t_1^2
		-(1-\omega_1-\omega_2)t_2^2,
		\label{eq:3}\\
		3\beta x^2
		&=
		\omega_1+\omega_2t_1^3
		-(1-\omega_1-\omega_2)t_2^3.
		\label{eq:4}
	\end{align}
	
	\subsection*{3. Нахождение $\omega_2$}
	
	Из \eqref{eq:1}:
	
	\[
	0
	=
	-\omega_1+\omega_2-1+\omega_1+\omega_2.
	\]
	
	Следовательно,
	
	\[
	0=2\omega_2-1,
	\]
	
	поэтому
	
	\begin{equation}
		\boxed{\omega_2=\frac12}.
	\end{equation}
	
	Тогда
	
	\[
	1-\omega_1-\omega_2
	=
	1-\omega_1-\frac12
	=
	\frac12-\omega_1.
	\]
	
	Таким образом,
	
	\begin{equation}
		\boxed{
			1-\omega_1-\omega_2
			=
			\frac12-\omega_1.
		}
	\end{equation}
	
	\subsection*{4. Второе скалярное уравнение}
	
	Подставляем найденные выражения в \eqref{eq:2}:
	
	\[
	\beta
	=
	\omega_1+\frac{t_1}{2}
	-
	\left(\frac12-\omega_1\right)t_2.
	\]
	
	Используем
	
	\[
	t_2=2t_1+1.
	\]
	
	Тогда
	
	\[
	\beta
	=
	\omega_1+\frac{t_1}{2}
	-
	\left(\frac12-\omega_1\right)(2t_1+1).
	\]
	
	Раскрываем скобки:
	
	\[
	\beta
	=
	\omega_1+\frac{t_1}{2}
	-\frac{2t_1+1}{2}
	+\omega_1(2t_1+1).
	\]
	
	Получаем
	
	\[
	\beta
	=
	\omega_1+\frac{t_1}{2}
	-t_1-\frac12
	+2\omega_1t_1+\omega_1.
	\]
	
	Следовательно,
	
	\[
	\beta
	=
	2\omega_1t_1+2\omega_1
	-\frac{t_1}{2}-\frac12.
	\]
	
	Вынося общий множитель, получаем
	
	\begin{equation}
		\boxed{
			\beta
			=
			(t_1+1)
			\left(
			2\omega_1-\frac12
			\right).
		}
		\label{eq:beta}
	\end{equation}
	
	\subsection*{5. Третье скалярное уравнение}
	
	Теперь рассмотрим \eqref{eq:3}:
	
	\[
	2\beta x
	=
	-\omega_1
	+\frac{t_1^2}{2}
	-
	\left(\frac12-\omega_1\right)t_2^2.
	\]
	
	Поскольку
	
	\[
	t_2=2t_1+1,
	\]
	
	то
	
	\[
	t_2^2=(2t_1+1)^2
	=
	4t_1^2+4t_1+1.
	\]
	
	Следовательно,
	
	\[
	2\beta x
	=
	-\omega_1
	+\frac{t_1^2}{2}
	-
	\left(\frac12-\omega_1\right)
	(4t_1^2+4t_1+1).
	\]
	
	Раскрываем скобки:
	
	\begin{align*}
		2\beta x
		={}&
		-\omega_1+\frac{t_1^2}{2}
		-\frac12(4t_1^2+4t_1+1)
		\\
		&+
		\omega_1(4t_1^2+4t_1+1).
	\end{align*}
	
	То есть
	
	\begin{align*}
		2\beta x
		={}&
		-\omega_1+\frac{t_1^2}{2}
		-2t_1^2-2t_1-\frac12
		\\
		&+
		4\omega_1t_1^2+4\omega_1t_1+\omega_1.
	\end{align*}
	
	Слагаемые $-\omega_1$ и $+\omega_1$ сокращаются:
	
	\[
	2\beta x
	=
	\frac{t_1^2}{2}
	-2t_1^2
	-2t_1
	-\frac12
	+
	4\omega_1t_1^2
	+
	4\omega_1t_1.
	\]
	
	Таким образом,
	
	\begin{equation}
		\boxed{
			2\beta x
			=
			4\omega_1t_1(t_1+1)
			-\frac32t_1^2
			-2t_1
			-\frac12.
		}
		\label{eq:third_correct}
	\end{equation}
	
	Теперь из \eqref{eq:beta}
	
	\[
	2\beta x
	=
	2x(t_1+1)
	\left(
	2\omega_1-\frac12
	\right),
	\]
	
	то есть
	
	\[
	2\beta x
	=
	x(t_1+1)(4\omega_1-1).
	\]
	
	Приравнивая это выражение к \eqref{eq:third_correct}, получаем
	
	\[
	x(t_1+1)(4\omega_1-1)
	=
	4\omega_1t_1(t_1+1)
	-\frac32t_1^2
	-2t_1
	-\frac12.
	\]
	
	Переносим члены с $\omega_1$ в левую часть:
	
	\[
	4\omega_1x(t_1+1)
	-
	4\omega_1t_1(t_1+1)
	=
	x(t_1+1)
	-\frac32t_1^2
	-2t_1
	-\frac12.
	\]
	
	Выносим $4\omega_1(t_1+1)$:
	
	\[
	4\omega_1(t_1+1)(x-t_1)
	=
	x(t_1+1)
	-\frac32t_1^2
	-2t_1
	-\frac12.
	\]
	
	Следовательно,
	
	\begin{equation}
		\boxed{
			\omega_1
			=
			\frac{
				2x(t_1+1)-3t_1^2-4t_1-1
			}{
				8(t_1+1)(x-t_1)
			}.
		}
		\label{eq:w1_intermediate}
	\end{equation}
	
	\subsection*{6. Нахождение $\beta$}
	
	Используем \eqref{eq:beta}:
	
	\[
	\beta
	=
	(t_1+1)
	\left(
	2\omega_1-\frac12
	\right).
	\]
	
	Подставляем \eqref{eq:w1_intermediate}:
	
	\[
	\beta
	=
	(t_1+1)
	\left[
	\frac{
		2x(t_1+1)-3t_1^2-4t_1-1
	}{
		4(t_1+1)(x-t_1)
	}
	-\frac12
	\right].
	\]
	
	Приводим к общему знаменателю:
	
	\[
	\beta
	=
	(t_1+1)
	\frac{
		2x(t_1+1)-3t_1^2-4t_1-1
		-2(t_1+1)(x-t_1)
	}{
		4(t_1+1)(x-t_1)
	}.
	\]
	
	Раскрываем последнюю скобку:
	
	\[
	2(t_1+1)(x-t_1)
	=
	2x(t_1+1)-2t_1(t_1+1).
	\]
	
	Поэтому числитель равен
	
	\begin{align*}
		&
		2x(t_1+1)-3t_1^2-4t_1-1
		\\
		&\qquad
		-2x(t_1+1)+2t_1(t_1+1)
		\\
		={}&
		-3t_1^2-4t_1-1
		+2t_1^2+2t_1
		\\
		={}&
		-t_1^2-2t_1-1
		\\
		={}&
		-(t_1+1)^2.
	\end{align*}
	
	Следовательно,
	
	\[
	\beta
	=
	-\frac{(t_1+1)^2}{4(x-t_1)}.
	\]
	
	Итак,
	
	\begin{equation}
		\boxed{
			\beta
			=
			\frac{(t_1+1)^2}{4(t_1-x)}.
		}
		\label{eq:beta_final}
	\end{equation}
	
	\subsection*{7. Четвёртое скалярное уравнение}
	
	Остаётся использовать \eqref{eq:4}:
	
	\[
	3\beta x^2
	=
	\omega_1+\frac12t_1^3
	-
	\left(\frac12-\omega_1\right)t_2^3.
	\]
	
	Так как
	
	\[
	t_2=2t_1+1,
	\]
	
	имеем
	
	\[
	t_2^3
	=
	(2t_1+1)^3
	=
	8t_1^3+12t_1^2+6t_1+1.
	\]
	
	Поэтому
	
	\[
	3\beta x^2
	=
	\omega_1+\frac12t_1^3
	-
	\left(\frac12-\omega_1\right)
	(8t_1^3+12t_1^2+6t_1+1).
	\]
	
	Подставляем \eqref{eq:w1_intermediate} и \eqref{eq:beta_final}. После раскрытия скобок и приведения подобных членов получаем
	
	\begin{equation}
		(t_1+1)^2
		\left(
		2t_1^2-6xt_1-2t_1+3x^2-1
		\right)=0.
	\end{equation}
	
	Поскольку $t_1\neq -1$, множитель $(t_1+1)^2$ ненулевой. Поэтому
	
	\[
	2t_1^2-6xt_1-2t_1+3x^2-1=0.
	\]
	
	Или
	
	\begin{equation}
		\boxed{
			2t_1^2-(6x+2)t_1+3x^2-1=0.
		}
		\label{eq:quadratic}
	\end{equation}
	
	\subsection*{8. Нахождение $t_1$}
	
	Дискриминант уравнения \eqref{eq:quadratic} равен
	
	\begin{align*}
		D
		&=
		(6x+2)^2
		-
		8(3x^2-1)
		\\
		&=
		36x^2+24x+4
		-24x^2+8
		\\
		&=
		12x^2+24x+12
		\\
		&=
		12(x+1)^2.
	\end{align*}
	
	Поэтому
	
	\[
	t_1
	=
	\frac{6x+2\pm\sqrt{12(x+1)^2}}{4}.
	\]
	
	Удобно записать два корня через параметр
	
	\[
	s\in\{-1,1\}.
	\]
	
	Тогда
	
	\begin{equation}
		\boxed{
			t_1
			=
			\frac{1+3x+s\sqrt3(x+1)}{2}.
		}
		\label{eq:t1_final}
	\end{equation}
	
	\subsection*{9. Нахождение $t_2$ и $t_3$}
	
	Из \eqref{eq:t2}
	
	\[
	t_2=2t_1+1.
	\]
	
	Подставляя \eqref{eq:t1_final}, получаем
	
	\begin{align*}
		t_2
		&=
		1+3x+s\sqrt3(x+1)+1
		\\
		&=
		2+3x+s\sqrt3(x+1).
	\end{align*}
	
	Следовательно,
	
	\begin{equation}
		\boxed{
			t_2
			=
			2+3x+s\sqrt3(x+1).
		}
		\label{eq:t2_final}
	\end{equation}
	
	А из \eqref{eq:t3}
	
	\[
	t_3=3t_1+2.
	\]
	
	Поэтому
	
	\begin{align*}
		t_3
		&=
		\frac{3+9x+3s\sqrt3(x+1)}{2}+2
		\\
		&=
		\frac{
			7+9x+3s\sqrt3(x+1)
		}{2}.
	\end{align*}
	
	Таким образом,
	
	\begin{equation}
		\boxed{
			t_3
			=
			\frac{
				7+9x+3s\sqrt3(x+1)
			}{2}.
		}
		\label{eq:t3_final}
	\end{equation}
	
	\subsection*{10. Нахождение $\omega_1$}
	
	Возвращаемся к выражению \eqref{eq:w1_intermediate}:
	
	\[
	\omega_1
	=
	\frac{
		2x(t_1+1)-3t_1^2-4t_1-1
	}{
		8(t_1+1)(x-t_1)
	}.
	\]
	
	Подставляем
	
	\[
	t_1=
	\frac{1+3x+s\sqrt3(x+1)}2.
	\]
	
	Сначала найдём
	
	\[
	t_1+1
	=
	\frac{
		3+3x+s\sqrt3(x+1)
	}{2}
	=
	\frac{x+1}{2}(3+s\sqrt3),
	\]
	
	и
	
	\[
	t_1-x
	=
	\frac{
		1+x+s\sqrt3(x+1)
	}{2}
	=
	\frac{x+1}{2}(1+s\sqrt3).
	\]
	
	Используя также то, что $s^2=1$, после упрощения выражения для $\omega_1$ получаем
	
	\begin{equation}
		\boxed{
			\omega_1
			=
			\frac{2+s\sqrt3}{8}.
		}
		\label{eq:w1_final}
	\end{equation}
	
	Здесь знак перед $\sqrt3$ тот же самый, что и в формуле для $t_1$.
	
	\subsection*{11. Нахождение $\beta$}
	
	Используем формулу \eqref{eq:beta_final}:
	
	\[
	\beta
	=
	\frac{(t_1+1)^2}{4(t_1-x)}.
	\]
	
	Подставляем
	
	\[
	t_1+1
	=
	\frac{x+1}{2}(3+s\sqrt3)
	\]
	
	и
	
	\[
	t_1-x
	=
	\frac{x+1}{2}(1+s\sqrt3).
	\]
	
	Получаем
	
	\[
	\beta
	=
	\frac{
		\dfrac{(x+1)^2}{4}(3+s\sqrt3)^2
	}{
		2(x+1)(1+s\sqrt3)
	}.
	\]
	
	Следовательно,
	
	\[
	\beta
	=
	\frac{x+1}{8}
	\frac{(3+s\sqrt3)^2}{1+s\sqrt3}.
	\]
	
	Так как
	
	\[
	(3+s\sqrt3)^2
	=
	12+6s\sqrt3
	=
	6(2+s\sqrt3),
	\]
	
	то
	
	\[
	\beta
	=
	\frac{3(x+1)(2+s\sqrt3)}
	{4(1+s\sqrt3)}.
	\]
	
	Рационализируя знаменатель,
	
	\[
	\frac{2+s\sqrt3}{1+s\sqrt3}
	=
	1+s\sqrt3,
	\]
	
	поскольку
	
	\[
	(1+s\sqrt3)^2
	=
	4+2s\sqrt3
	=
	2(2+s\sqrt3).
	\]
	
	Поэтому
	
	\begin{equation}
		\boxed{
			\beta
			=
			\frac{3(1+s\sqrt3)(x+1)}{8}.
		}
		\label{eq:beta_final_s}
	\end{equation}
	
	\subsection*{12. Окончательный ответ}
	
	Итак, при
	
	\[
	s=\pm1
	\]
	
	получаем две ветви решения:
	
	\begin{equation}
		\boxed{
			\omega_1
			=
			\frac{2+s\sqrt3}{8},
		}
	\end{equation}
	
	\begin{equation}
		\boxed{
			\omega_2
			=
			\frac12,
		}
	\end{equation}
	
	\begin{equation}
		\boxed{
			t_1
			=
			\frac{1+3x+s\sqrt3(x+1)}{2},
		}
	\end{equation}
	
	\begin{equation}
		\boxed{
			t_2
			=
			2+3x+s\sqrt3(x+1),
		}
	\end{equation}
	
	\begin{equation}
		\boxed{
			t_3
			=
			\frac{7+9x+3s\sqrt3(x+1)}{2},
		}
	\end{equation}
	
	\begin{equation}
		\boxed{
			\beta
			=
			\frac{3(1+s\sqrt3)(x+1)}{8}.
		}
	\end{equation}
	
	Таким образом,
	
	\begin{equation}
		\boxed{
			(\omega_1,\omega_2,t_1,t_2,t_3,\beta)
			=
			\left(
			\frac{2+s\sqrt3}{8},
			\frac12,
			\frac{1+3x+s\sqrt3(x+1)}2,
			2+3x+s\sqrt3(x+1),
			\frac{7+9x+3s\sqrt3(x+1)}2,
			\frac{3(1+s\sqrt3)(x+1)}8
			\right),
		}
	\end{equation}
	
	где
	
	\[
	\boxed{s=\pm1}.
	\]
	
	При этом
	
	\[
	1-\omega_1-\omega_2
	=
	\frac{2-s\sqrt3}{8}.
	\]
	
	Действительно,
	
	\[
	t_2-t_1=t_1+1
	\]
	
	и
	
	\[
	t_3-t_2=t_1+1,
	\]
	
	поэтому оба дополнительных условия выполняются.
	
\end{document}